\scanpage{52}
\margin{51}
\mspage{-3.18-}

form $\bar{\Phi}_0^{-1}(V)$ for some open neighborhood $V$ of $y$. Since $U_0$ is affine, it follows (for example, by SGA 8, Cor. 5.6) that $V$ is affine. Moreover, $U_0$ is defined by $P\ne0$, where $P$ is an invariant section of $(L_0\otimes\cdots\otimes L_m)^N$ (for suitable $N$). Then $P^{n+1}$ is induced by a section $s$ of $(M_0\otimes\cdots\otimes M_m)^N$ (cf. \S1), and $V$ must be defined by $s\ne0$. \uline{QED}

For the sake of the sequel, we fix a more explicit notation:
$(\mathbb{P}_n)_{\mathrm{stable}}^{m+1}$ instead of $U_{\mathrm{stable}}$, (cf. Proposition 7.7).
